\section{Validazione dello Schema}
\label{sec:validazione}

\subsection{validate\_tbox.py}

Lo script \texttt{validate\_tbox.py} implementa un validatore statico dello schema
che esegue \textbf{10 controlli di consistenza} sullo schema dell'ontologia.
La Tabella~\ref{tab:checks} elenca i controlli con codice e severit\`a.

\begin{table}[h]
\centering
\caption{Controlli di validazione dello schema}
\label{tab:checks}
\begin{tabular}{cll}
\toprule
\textbf{\#} & \textbf{Codice}                 & \textbf{Descrizione} \\
\midrule
1 & DANGLING\_REFERENCE        & Riferimenti a entit\`a non dichiarate (ERROR) \\
2 & UNSATISFIABLE\_DISJOINT    & Classi con superclassi disgiunte (ERROR) \\
3 & CIRCULAR\_SUBCLASS         & Cicli nella gerarchia rdfs:subClassOf (ERROR) \\
4 & INVERSE\_CONSISTENCY       & Consistenza owl:inverseOf bidirezionale (WARNING) \\
5 & PROPERTY\_TYPE\_CONFLICT    & Tipi di propriet\`a mutuamente esclusivi (ERROR) \\
6 & MISSING\_DOMAIN\_RANGE      & Domain/range mancanti su propriet\`a (WARNING) \\
7 & MISSING\_LABEL             & Label rdfs:label mancanti (WARNING) \\
8 & UNUSED\_CLASSES            & Classi dichiarate mai referenziate (WARNING) \\
9 & RESTRICTION\_COMPLETENESS  & Restrizioni incomplete (ERROR/WARNING) \\
10& DATATYPE\_REST\_SYNTAX      & Sintassi datatype restriction (WARNING) \\
\bottomrule
\end{tabular}
\end{table}

\subsection{Risultato della Validazione}

La validazione sullo schema attuale (\texttt{videogames.owl}) produce il seguente
risultato:

\begin{itemize}[nosep]
  \item \textbf{0 ERROR} --- nessuna violazione critica;
  \item \textbf{18 WARNING} --- tutti benigni e giustificati;
  \item \textbf{Esito}: PASS.
\end{itemize}

\subsection{Analisi dei 18 Warning}

I 18 warning si dividono in due categorie:

\subsubsection{12 Warning INVERSE\_NOT\_BIDIRECTIONAL}

Il validatore segnala che, per 12 coppie inverse, solo una direzione dichiara
esplicitamente \texttt{owl:inverseOf} (es. \texttt{developedBy owl:inverseOf developerOf}
ma non viceversa). Questo \`e intenzionale:
\begin{itemize}[nosep]
  \item \texttt{owlrl} inferisce automaticamente la direzione opposta durante la
    materializzazione, quindi la dichiarazione bidirezionale sarebbe ridondante;
  \item Lo schema originale V1 dichiarava solo la direzione ``forward''; la V2 ha
    aggiunto le inverse come nuove propriet\`a, ma l'\texttt{inverseOf} inverso
    non \`e stato aggiunto per non duplicare assiomi gi\`a inferibili.
\end{itemize}

\subsubsection{6 Warning MISSING\_DOMAIN / MISSING\_RANGE}

Riguardano sei propriet\`a senza dominio o range dichiarato:
\begin{itemize}[nosep]
  \item \texttt{name}, \texttt{description} --- super-propriet\`a generiche
    intenzionalmente polimorfiche (usate come base per gameName, developerName,
    ecc. tramite \texttt{rdfs:subPropertyOf}), quindi senza dominio fisso;
  \item \texttt{hasWikidataId} --- propriet\`a trasversale applicabile a qualsiasi
    risorsa con un identificativo Wikidata, intenzionalmente senza dominio;
  \item \texttt{subsidiaryOf} --- senza dominio/range per consentire l'uso sia
    su Developer che su Publisher;
  \item \texttt{reviewer} --- range non specificato perch\'e il revisore pu\`o
    essere un'entit\`a testuale o un URI esterno.
\end{itemize}

Tutti questi warning sono giustificati da scelte di design e non costituiscono
problemi di qualit\`a dell'ontologia.
